% Eleventh Philippic --- headnote
% philippic-11, March 43 BC, Rome

\textit{Delivered in the Senate in the early spring of 43 BC, the
Eleventh Philippic answers a single piece of dreadful news. Gaius
Trebonius --- one of Caesar's assassins, now proconsul of Asia ---
had been surprised at night in Smyrna by Publius Cornelius
Dolabella, who had marched east through the province with a legion;
tortured for two days over the public money, killed, beheaded, and
his body flung into the sea. The day before, on the motion of
Quintus Fufius Calenus, the Senate had declared Dolabella a public
enemy and ordered his goods confiscated. The present debate is the
sequel: who is to be given the war against him. Cicero opens over
the body of Trebonius --- the cruelty of Dolabella made the measure
of what Antony, the partner of his madness, would do in Italy if he
could --- and argues from the atrocity that the war for liberty and
the war against Dolabella are one and the same.}

\textit{Two motions are already before the house, and Cicero rejects
both. The first, Lucius Caesar's, would confer an extraordinary
command on a distinguished private citizen, Publius Servilius
Isauricus; Cicero opposes it on the old principle that extraordinary
commands are ``a popular, windy thing,'' unworthy of the Senate's
gravity, and refuses to let the curia be turned into an electoral
assembly where senators canvass and a friend's defeat becomes a
personal slight. (His own grant to the young Octavian, he answers
the obvious objection, only ratified a command the facts had already
created: ``what was not snatched away is not to be reckoned as
given.'') The second motion would send the consuls themselves,
Pansa and Hirtius, to draw lots for Asia and Syria once Decimus
Brutus is freed --- which Cicero calls shameful and useless while a
consul-elect is besieged at Mutina and the whole war turns on his
rescue. From this comes the speech's most quoted sentence
(\textsection 28), the natural-law defence of Cassius's having
acted without orders: law is nothing else than right reason, drawn
from the will of the gods; and a province seized to save the
commonwealth is held, the written statutes once trampled down, ``by
the law of nature.''}

\textit{Cicero's own proposal, set out as a formal decree
(\textsection 29--31), confirms Gaius Cassius --- already in the
East to keep Dolabella out of Syria --- in command of that province
``by the best right,'' to take over the armies of Marcius Crispus,
Staius Murcus, and Allienus, to requisition ships, men, and money
across Syria, Asia, Bithynia, and Pontus, and to pursue Dolabella by
land and sea, with the elder and younger Deiotarus and the other
client kings invited to aid him. The case for Cassius is built from
his proven army, his standing among the Tyrians and in Phoenicia,
and his old victory over the Parthians --- his greatest title to
honour, the part in Caesar's death, being pointedly left to the
``testimony of memory rather than of speech.'' Around the argument
runs the familiar Philippic invective, a contemptuous muster of
Antony's creatures (his brother Lucius, the bankrupt Trebellius
``turned into a new account-book himself,'' Plancus, the
land-commissioners Nucula and Lento, Annius Cimber, the squatters
Saxa and Cafo), and the recurring retort to those who plead the
veterans' feelings: the Senate will not choose its commanders ``at
the veterans' nod.'' In the event the motion did not carry --- the
house preferred to leave the command to the consuls --- but Cassius
made the speech prophetic, taking Syria, raising the legions Cicero
named, and driving Dolabella to suicide at Laodicea before the
summer was out.}
